% Source: https://artofproblemsolving.com/community/c3h1584717_systems_of_literal_equations
% Stable identity: aops-topic-1584717-q01
From the second equation $x=y+2b$, so with $Y=y+1$ (which must be nonzero) the first equation becomes
\[
\dfrac{Y+2b}{Y}=\dfrac{a+b+1}{a-b+1}.
\]
Cross-multiplying,
\[
(Y+2b)(a-b+1)=Y(a+b+1)\qquad\Longleftrightarrow\qquad 2b(a-b+1)=Y\left[(a+b+1)-(a-b+1)\right]=2bY .
\]
Since $b\ne0$, $Y=a-b+1$, which is nonzero by assumption. Hence $y=a-b$ and $x=y+2b=a+b$:
\[
\boxed{x=a+b,\quad y=a-b}.
\]
Check: $\dfrac{x+1}{y+1}=\dfrac{a+b+1}{a-b+1}$ and $x-y=2b$.

(If $b=0$, the equations only say $x=y$ with $y\ne-1$, so the solution is not unique; this is why $b\ne0$ is assumed.)
